% Feature engineering (issue #16), based on the completed notebook in issue #15.
% Citation keys are shared with the literature review; BibTeX is supplied in #22
% for consolidation into reports/phase1/refs.bib.

\section{Feature Engineering}
\label{sec:feature_engineering}

Phase~1 represents each engine cycle by a fixed-length vector of operating
context, sensor levels and causal summaries of recent history. This follows the
separation of health-indicator construction from RUL prediction discussed by
Lei et al.~\cite{Lei2018machinery}. Wang et al.~\cite{Wang2008similarity}
estimate remaining life by matching degradation patterns to historical engines;
our regressors instead receive explicit summaries of those patterns. The
definitions below describe the implementation in \texttt{rul.features} and the
completed feature-engineering notebook.

\subsection{Sensor Selection and Normalized Readings}
Sensor selection uses raw training-fold rows before normalization. The default
rule removes only exactly constant channels; it does not assume that every
low-variance sensor is uninformative. An optional relative-standard-deviation
cutoff compares the population standard deviation with
$\max(|\mathrm{mean}|,1)$ and must be chosen using training data only. On all
official FD001 training engines, the default rule removes sensors
$s_1,s_5,s_{10},s_{16},s_{18},s_{19}$, retaining 15 channels, including
$s_6$, which is nearly but not exactly constant (two distinct values). The
retained list is reused for held-out engines. This reduces
redundant inputs without supervised feature selection. Khelif
et al.~\cite{Khelif2017direct} motivate selecting informative sensors for direct
RUL regression, although their wrapper selection differs from our variability
filter.

Operating conditions can shift sensor levels independently of degradation.
K-means is fitted to the three unscaled operating settings from training rows,
with one cluster for FD001/FD003 and six for FD002/FD004. The implementation
uses seed~42 and ten initializations. Each reading is assigned to its nearest
learned centroid, giving regime $r$, and transformed as
\begin{equation}
    z_{u,t,s}=\frac{x_{u,t,s}-\mu_{r,s}}{\sigma_{r,s}},
\end{equation}
where $u,t,s$ denote engine, cycle and sensor, and the mean and population
standard deviation are fitted within training regime $r$. A zero standard
deviation uses a scale of one. Retained instantaneous $z_{u,t,s}$ values remain
predictors, preserving current condition information that averaging may hide.
The FD002 training-only comparison for $s_2,s_3,s_{11}$ reduces regime-mean
dispersion from about one global standard deviation to numerical zero. This
checks removal of training mean offsets; it does not establish a prediction
gain.

\subsection{Cycle and Operating Context}
The current \texttt{cycle} supplies an observed age proxy, allowing models to
relate accumulated operation to remaining life. It contains no final-lifetime
information, so engines of the same age can still receive different predictions
from their sensors. The original \texttt{op\_1}, \texttt{op\_2},
\texttt{op\_3} settings and learned \texttt{regime} are also retained to describe
the current operating context after sensor normalization. Regime is stored as
an integer in the matrix; cluster numbers are arbitrary identifiers with no
physical ordering. FD001 has a constant regime column, retained for a
consistent schema.

\subsection{Rolling Degradation Features}
For each retained normalized sensor, trailing windows of $w\in\{5,10,20\}$
observations summarize short and longer histories. Let $\mathcal{H}_{u,t,w}$
contain the current and up to $w-1$ preceding observations of engine $u$, and
let $n=|\mathcal{H}_{u,t,w}|$. Rows are sorted by engine and cycle for
calculation, then restored to their input order. Windows never cross engines,
use no future readings, and use the available history at startup.

\subsubsection{Rolling mean}
The local level is
\begin{equation}
    m_{u,t,s,w}=\frac{1}{n}\sum_{i\in\mathcal{H}_{u,t,w}}z_{u,i,s}.
\end{equation}
Averaging smooths cycle-to-cycle fluctuations. Short windows respond faster;
long windows smooth more strongly but can delay changes. These window sizes
are implementation defaults, not claimed optimal values.

\subsubsection{Rolling slope}
The least-squares trend over the same history is
\begin{equation}
    b_{u,t,s,w}=
    \frac{\sum_{i\in\mathcal{H}_{u,t,w}}(c_i-\bar{c})(z_{u,i,s}-m_{u,t,s,w})}
         {\sum_{i\in\mathcal{H}_{u,t,w}}(c_i-\bar{c})^2},
\end{equation}
where $c_i$ is the recorded cycle and $\bar{c}$ its window mean. The slope is
zero for a single observation. It captures direction and rate of change in
normalized sensor units per cycle, complementing the local level; short-window
slopes can be sensitive to noise.

\subsubsection{Signed health change}
The engine-relative change is
\begin{equation}
    h_{u,t,s,w}=m_{u,t,s,w}-z_{u,t_0,s},
\end{equation}
where $t_0$ is that engine's first observed cycle. This provides a baseline
relative to each engine's initial observed condition. Unlike a calibrated
health index, it need not be monotonic or lie in $[0,1]$; the first observation
is not assumed fault-free. It is zero initially, and its sign depends on the
sensor's degradation direction. Rolling summaries use regime-normalized
readings and may span regime changes within an engine; they do not restart at
each change of regime.

\subsection{Feature Matrix and Leakage Control}
With $S$ retained sensors and $W$ window sizes, the predictor count is
$5+S(1+3W)$: cycle, three settings, regime, $S$ instantaneous readings and
$3SW$ rolling features. The FD001 export therefore has 155 predictors for
$S=15$, $W=3$, with 20,631 training rows and 13,096 test rows across 100 engines
in each split. The notebook writes aligned train/test CSVs and a JSON manifest
under \texttt{data/processed/}; the manifest records sensors, predictor order,
windows, regime count, seed and fitting scope. Engine identifier
\texttt{unit} is metadata, and uncapped training RUL is a target; neither is a
predictor. Test features contain no target.

Final exports fit selection and normalization on all official training engines.
For engine-grouped cross-validation, both must instead be refitted within each
training fold; globally exported features cannot be reused for tuning. Rolling
features are computed independently within each split using only observed
engine history. Official test scoring uses each engine's last observed cycle
and the corresponding NASA RUL label. The notebook's training correlations
support exploratory interpretation only: the strongest is \texttt{cycle}
itself ($\rho=-0.79$), several rolling means have absolute Spearman
correlations near 0.75 with RUL, and slopes are noisier. These
observations do not establish held-out accuracy or a benefit from any individual
feature group; that requires model evaluation and ablations.
